\subsection{由半范数定义的拓扑}

设 \(p\) 是向量空间 \(E\)（在 \(K\) 上）上的一个半范数；对每个 \(\alpha > 0\)，用 \(V(p, \alpha)\) 表示由那些 \(x\)（属于 \(E\) 且满足 \(p(x) \leqslant \alpha\)）组成的子集。显然，若 \(x \in V(p, \alpha)\) 且 \(\lambda \in K\) 满足 \(|\lambda| \leqslant 1\)，则 \(\lambda x \in V(p, \alpha)\)，换言之，\(V(p, \alpha)\) 是均衡的。又，对每个 \(x_0 \in E\)，存在非零标量 \(\mu \in K\) 使得 \(|\mu| \geqslant p(x_0) \alpha^{-1}\)，因此 \(\mu^{-1}x_0 \in V(p, \alpha)\)，也就是说 \(V(p, \alpha)\) 是吸收的。最后，由 \((\mathrm{SN}_{\mathrm{II}})\) 有 \(V(p, \alpha/2) + V(p, \alpha/2) \subset V(p, \alpha)\)，并由 \((\mathrm{SN}_{\mathrm{I}})\)，对每个非零标量 \(\lambda\)（在 \(K\) 中）有 \(\lambda V(p, \alpha) = V(p, |\lambda| \alpha)\)。由这些备注并根据 I, 第 7 页 的 prop. 4，我们断定：当 \(\alpha\) 在 \(> 0\) 的数集中变化（或只在趋于 0 的一个严格正数序列中变化）时，诸集 \(V(p, \alpha)\) 构成 \(E\) 的一个与向量空间结构相容的拓扑的 0 的基本邻域系；我们称这拓扑由半范数 \(p\) 定义。赋有这种拓扑的向量空间 \(E\) 称为半赋范空间。注意，若 \(W(p, \alpha)\) 表示由那些 \(x\)（属于 \(E\) 且满足 \(p(x) < \alpha\)）组成的子集，则诸 \(W(p, \alpha)\)（其中 \(\alpha > 0\)，或 \(\alpha\)）

在一个趋于零的严格正数序列中变化）构成由 \(p\) 定义的拓扑的 0 的一个基本邻域系。

若 \(\Gamma\) 是 \(\mathbf{E}\) 上半范数的一个集合，则由诸半范数 \(p \in \Gamma\) 定义的拓扑的上确界与向量空间结构相容（I, 第 10 页, cor. 4）。这个拓扑的 0 的一个基本邻域系由诸有限交 \(\bigcap_{i} \mathrm{V}(p_i, \alpha_i)\) 给出，其中 \(p_i \in \Gamma\) 且 \(\alpha_i > 0\)。称这拓扑是由半范数集

\(\Gamma\) 定义的。它是 \(\mathbf{E}\) 上在所有平移下不变且使诸半范数 \(p \in \Gamma\) 连续的拓扑中最粗的一个。

设 E 是 K 上的一个拓扑向量空间：E 上的半范数组 \(\Gamma\) 称为基本半范数系，如果 E 上的拓扑与由 \(\Gamma\) 定义的拓扑相同。

设 E 是 K 上的向量空间，赋有由半范数集 \(\Gamma\) 定义的拓扑。对每个半范数 p，有 \(p(x - z) \leqslant p(x - y) + p(y - z)\)，这表明函数 \((x, y) \mapsto p(x - y)\) 是 E 上的一个伪距离 (GT, IX, §1.1)：由定义可得，当 p 在 \(\Gamma\) 中变化时，这些伪距离组成的集定义了拓扑向量空间 E 的一致结构。

注记. --- 1) E 上由有限个半范数 \(p_i\)（\(1 \leqslant i \leqslant n\)）定义的拓扑可由单个半范数 \(p = \sup_{1 \leqslant i \leqslant n} p_i\) 定义。但由

无穷多个半范数定义的拓扑一般不能由单个半范数定义（III, 第 37 页, exerc. 2）。

\begin{enumerate}
\def\labelenumi{\arabic{enumi})}
\setcounter{enumi}{1}
\item
  设 \((\mathcal{T}_{\iota})_{\iota\in I}\) 是 K 上向量空间 E 上的一族拓扑，其中每个拓扑都由一族半范数 \(\Gamma_{\iota}\) 定义。则由半范数集 \(\Gamma = \bigcup \Gamma_{\iota}\) 定义的拓扑是诸拓扑 \(T_{\iota}\) 的上确界。
\item
  若 \(\Gamma_0\) 是一个半范数集，它由 E 上两个半范数 \(p, q\) 之间如下定义的递增序关系所定向：「存在 \(\lambda > 0\) 使得 \(p \leqslant \lambda q\)」，则对由 \(\Gamma_0\) 定义的拓扑，取诸集 \(\mathbf{V}(p, \alpha)\)（其中 \(p \in \Gamma_0\)，\(\alpha > 0\)）即得 0 的一个基本邻域系。若 \(\Gamma\) 是 E 上任意一个半范数集，则与 \(\Gamma\) 定义相同拓扑的一个滤过化半范数集，就是集 \(\Gamma_0\)（由属于 \(\Gamma\) 的半范数的一切有限族的上包络所组成）。
\item
  即使 \(\mathbf{K} = \mathbf{R}\)，\(\mathbf{K}\) 上拓扑向量空间的拓扑也不总能由半范数集定义（cf.~II, 第 24 页, corollary）。
\end{enumerate}

例. --- 设 \(\mathcal{C}^{\infty}(\mathbf{R})\) 是 R 上在 R 中无穷次可微的实值函数组成的 R 上的向量空间。对每个函数以及每对整数 \(n \geqslant 0\)、\(m \geqslant 1\)，置

\[
p _ {n, m} (f) = \sup _ {- m \leqslant t \leqslant m} | f ^ {(n)} (t) |\tag{2}
\]

其中 \(f^{(0)} = f\)。显然诸 \(p_{n,m}\) 是 \(\mathcal{C}^{\infty}(\mathbf{R})\) 上的半范数。为使诸函数 \(f_{\alpha}\) 趋于 0（按指标集上的一个滤子 \(\mathfrak{F}\)），在 \(\mathcal{C}^{\infty}(\mathbf{R})\) 中关于拓扑 \(\mathcal{T}\)（由半范数 \(p_{n,m}\) 定义的），必须且只需：对所有整数 \(n \geqslant 0\)，诸函数 \(f_{\alpha}^{(n)}\)（按 \(\mathfrak{F}\)）在 \(\mathbf{R}\) 的每个紧子集上一致趋于 0。我们称 \(\mathcal{T}\) 是函数 \(f \in \mathcal{C}^{\infty}(\mathbf{R})\) 及其所有导数的紧收敛拓扑（cf.~III, 第 9 页）。

命题 2. --- 在向量空间 E 上，设 T 是由半范数集 Γ 定义的拓扑。

\begin{enumerate}
\def\labelenumi{(\roman{enumi})}
\item
  \(\{0\}\) 在 \(\mathbf{E}\) 中关于 \(\mathcal{T}\) 的闭包，是由诸 \(x \in \mathbf{E}\)（满足 \(p(x) = 0\)，对每个半范数 \(p \in \Gamma\)）组成的子集。
\item
  若 \(\mathcal{T}\) 是 Hausdorff 的且 \(\Gamma\) 可数，则 \(\mathcal{T}\) 可度量化。
\end{enumerate}

本命题由定义及 GT, IX, § 2.4, cor. 1 立即得出。

注意，即使 \(\mathcal{T}\) 可度量化，\(\mathcal{T}\) 也可能不能由单个范数来定义；上面给出的例子正是这种情形（cf.~IV, 第 18 页, 例 4）。

设 E 是 K 上的向量空间，赋有由半范数集 \(\Gamma\) 定义的拓扑。设 \(\hat{E}\) 是 E 的 Hausdorff 完备化（I, 第 6 页），\(\hat{\Gamma}\) 是映射 \(\hat{p}\)（\(\hat{E}\) 到 \(R_{+}\) 中的）组成的集，其中 p 在 \(\Gamma\) 中变化（GT, II, § 3.7, prop. 15）。由不等式延拓原理，诸函数 \(\hat{p} \in \hat{\Gamma}\) 是 \(\hat{E}\) 上的半范数，且诸函数 \(\hat{p}(x - y)\) 组成一个定义 \(\hat{E}\) 的一致结构的伪距离集（GT, IX, § 1.3, prop. 1）。因此我们看到，\(\hat{\Gamma}\) 是定义 \(\hat{E}\) 的拓扑的一个基本半范数集。

